\section{Case Studies}\label{sec:cases}

The four 2025--26 cases each isolate one mechanism of the model. Read together with three historical cases, they show how the mechanisms combine.

\subsection{Quinn Hughes: informal protection without a clause}

Hughes held no trade protection: his contract, which ran through 2026--27 at \$7.85 million, contained no clause \citep{cbs_hughes}. His agent says there was no formal trade request. What Vancouver had was a signal, given in the previous offseason, that a long-term extension was highly unlikely, and the trade followed from it \citep{espn_hughes}. Yet the agent described Vancouver's conduct this way: the club worked with the player ``as if we had a no-trade clause, almost'' \citep{dailyhive_hughes}. This is Proposition~\ref{prop:informal} in action. With eighteen months left on his deal, most of Hughes's value to an acquirer lay in the right to extend him, and that right was worth something only to teams he would agree to re-sign with. Selling to anyone else would have been selling a rental. The return was large: two young NHL roster players (Marco Rossi and Zeev Buium), a prospect (Liam \"Ohgren) and a first-round pick. That fits the model's claim that informal protection narrows the effective market without eliminating competition among the approved buyers. The postscript fits too. As of late October 2026, Hughes had not extended with Minnesota and was reported to prefer a short term \citep{rmnb_hughes}. The extension option Minnesota bought remains unexercised.

\subsection{Brady Tkachuk: the veto player who still got a price}

Tkachuk held a full NMC with two seasons left at \$8.2 million. Accounts differ on how his exit began. On the night of the trade, one outlet reported that ``a formal trade request may not have been made,'' only an unwillingness to extend \citep{fourthperiod_tkachuk}. The next day general manager Steve Staios said that ``what changed was a trade request,'' and Tkachuk confirmed he had asked out: ``It wasn't an easy decision'' \citep{yahoo_staios,thescore_tkachuk}. What remains disputed is when the decision was made, since Tkachuk had publicly called the rumours ``just not true'' in April. Staios was explicit about the clause's effect: ``The dynamic of a no-trade clause changes how the transaction works and how it's going to look\ldots It's not like he was available to every team in the league.'' The process was ``clearly pointed at one team at the end'' \citep{yahoo_staios,tsn_staios}. Tkachuk's reported list had four teams: Florida, Minnesota, Vegas and Carolina.

The model predicts a poor return when $m$ is effectively one, yet Ottawa received the ninth and 25th picks in the 2026 draft, a top-ten-protected 2029 first and a 2027 second \citep{fourthperiod_tkachuk}. That is the largest pick package in the dataset. Two features of the model reconcile the outcome. First, Ottawa's retention value was high. Tkachuk was under contract for two more seasons, the request was private until after the trade, and Ottawa plausibly could have kept him into the season. Under the outside-option logic of Section~\ref{sec:model}, a credible threat to keep the player sustains the price even in bilateral bargaining. Second, Florida's valuation was unusually high, because it was the only approved team that could pair Brady with his brother Matthew. A single buyer with a very high $v$ can pay a lot even when $T$'s share of the surplus is modest. The Tkachuk case is therefore not evidence against the veto mechanism. It is evidence that the mechanism operates on the seller's \emph{share}, and that the share can be protected by keeping the request private.

\subsection{Dylan Larkin: stalemate and sunk costs}

Larkin asked Detroit to trade him on 4 June 2026, with five seasons remaining at \$8.7 million on an eight-year contract he had signed as captain in 2023 \citep{wilx_larkin,nhl_larkin}. His initial list had three teams: Minnesota, Florida and Vegas. In the model's terms $m=3$, the region of Table~\ref{tab:model} where the seller loses roughly half the competitive price. Proposition~\ref{prop:delay} predicts delay in exactly these conditions, because a narrow list lowers the value of selling now relative to waiting. Delay is what followed.

Detroit's general manager, Steve Yzerman, was removed on 15 July. Shawn Horcoff was named interim general manager on 15 September \citep{nhl_horcoff}. Larkin was stripped of the captaincy, booed at the home opener, and placed on non-roster injured reserve, with the coach remarking that ``he's got to live with the ramifications of asking for a trade'' \citep{espn_larkin_ir}. The list expanded under pressure, to ten teams by late October \citep{dfo_larkin_list}. In the model this is $m$ rising over time, and it is the player, not the clause, that relaxes the constraint. On 29 October, after Larkin's agent said he would not report, Detroit suspended him without pay. Horcoff said the club would ``continue to explore trade opportunities,'' but only in ``the best interests of the Detroit Red Wings'' \citep{cod_larkin_suspended}. Suspension is the team's counter-move. It cannot remove the veto, but it raises the player's cost of waiting, which pushes him to widen the list.

Larkin is also the cleanest example of the sunk-cost temptation in Proposition~\ref{prop:delay}(c). A team that has committed \$69.6 million to its captain and built its rebuild around him has every reason to believe the relationship can be repaired, that is, to overestimate $\rho$. Detroit's first moves (waiting, then a leadership change) look like a bet on reconciliation. Its later moves (stripping the captaincy, the suspension) look like the recognition that $\rho$ had fallen to zero. At the time of writing the spell is censored at 148 days and counting.

\subsection{Connor Hellebuyck: audience costs}

Hellebuyck held an NMC with five seasons remaining at \$8.5 million. He asked for a trade privately early in the 2026 offseason. On 27 August he made the request public in a statement posted by his agent \citep{tsn_hellebuyck_open}. When he did not report to camp, Winnipeg suspended him on 16 September. General manager Kevin Cheveldayoff insisted that ``a contract is a contract'' \citep{nhl_hellebuyck_suspended}. Buffalo pursued him for weeks without a deal. On 23 October Winnipeg traded him to New Jersey for forward Dawson Mercer, the veteran Nick Bjugstad, goaltender Nico Daws, the rights to prospect Alexander Command, a 2029 second-round pick and a 2027 seventh \citep{nhl_hellebuyck_trade,tsn_hellebuyck_trade}.

The sequence follows the logic of equation~\eqref{eq:public}. While the request was private, Winnipeg could plausibly tell suitors it was prepared to keep the reigning Olympic best goaltender, so $R$ was high and selling was not urgent. Going public, then refusing to report, lowered $\lambda$ to zero and with it $R$. That made the exit credible and the sale inevitable. Unlike Larkin, Hellebuyck was reportedly open to a broad range of destinations \citep{thw_larkin_hellebuyck,tsn_hellebuyck_open}. Proposition~\ref{prop:audience} says the cost of publicity is smallest when $m$ is large, and the return was a useful but unspectacular package of roster players with no first-round pick. That is about what the model predicts for a seller with no outside option and a moderately competitive market.

\subsection{Historical cases}

\paragraph{Dany Heatley (2009): the veto exercised.} Heatley asked Ottawa to trade him in June 2009, with five years left on a six-year, \$45 million contract. When Ottawa agreed a deal with Edmonton for Andrew Cogliano, Dustin Penner and Ladislav \v{S}m\'{i}d, he refused to waive \citep{heatley_thn}. Ottawa settled for Milan Michalek, Jonathan Cheechoo and a second-round pick from San Jose in September. The case shows that the veto does more than shrink the market ex ante. It also lets the player reject a completed negotiation, which hands the remaining buyers knowledge of the seller's weakened position.

\paragraph{Martin St. Louis (2014): one destination.} St. Louis asked Tampa Bay privately in January 2014 and would accept only the Rangers \citep{stlouis_cbs}. With $m=1$, Tampa Bay's return was Ryan Callahan, a pending free agent, plus first-round picks. This is the bilateral case of the model, with a seller whose retention value of an aging star at the deadline was modest.

\paragraph{Roberto Luongo (2012--14): recapture.} Luongo's twelve-year contract and NTC, combined with a public willingness to move only to Florida, produced a stalemate of almost two years \citep{luongo_nesn}. Vancouver eventually moved him to Florida for a young goaltender, Jacob Markstr\"om, and a forward, Shawn Matthias. The long contract made the recapture liability enormous, which deterred buyers. The narrow list removed competition, and Vancouver's preference to wait for a better market cost it the asset's value. This is the longest-duration resolved case in the data and the clearest historical analogue for Larkin.

\paragraph{The counterfactuals.} Matt Duchene and Jack Eichel, the two most prominent unprotected stars who wanted out, also produced long standoffs (315 and 178 days). But their teams eventually ran open, competitive auctions and extracted large returns \citep{eichel_ap}. Duchene's three-team deal yielded Samuel Girard, prospects, and an Ottawa first that became the fourth overall pick. The absence of a veto did not prevent delay. It did preserve competition.
